% 引言章节

\section{引言}

\subsection{研究背景}

随着数字支付的普及，信用卡欺诈已成为全球金融体系面临的最严峻威胁之一。根据Nilson Report的最新数据，2023年全球信用卡欺诈损失预计超过350亿美元，且呈逐年增长趋势。传统的基于规则的风控系统在面对日益复杂的欺诈手段时显得力不从心——攻击者不断调整策略以规避固定规则，导致误报率居高不下，大量合法交易被误拦截，严重影响用户体验。

近年来，机器学习方法在欺诈检测领域展现出显著优势。相较于规则系统，机器学习模型能够自动从历史交易数据中学习欺诈模式，适应新型攻击手段。然而，欺诈检测面临一个根本性挑战：类别极端不平衡。在真实场景中，欺诈交易通常仅占总交易的0.1\%至1\%，这使得标准分类模型倾向于将所有交易预测为"正常"，从而导致欺诈漏检。

深度神经网络在计算机视觉和自然语言处理领域的成功引发了其在表格数据上的探索热潮。传统观点认为，梯度提升树（如XGBoost）在表格数据上的表现优于深度学习方法。但最新研究表明，专门为表格数据设计的深度架构——如FT-Transformer和DeepFM——在特定场景下能够接近甚至超越树模型。然而，深度模型在表格数据上的核心瓶颈并非模型容量，而是归纳偏置的错配：树模型的轴对齐分裂天然契合表格特征的异构性，而深度模型的全连接层和全局注意力缺乏这一先验。

\subsection{研究问题}

本文旨在回答以下核心研究问题：

\begin{enumerate}
    \item 深度表格模型（FT-Transformer、DeepFM）在信用卡欺诈检测任务上的表现是否能接近或超越XGBoost基线？自注意力机制和显式特征交互两种不同的归纳偏置在PCA匿名化特征空间中表现如何？
    \item 能否通过知识蒸馏将XGBoost的树归纳偏置（软标签分布 + 特征分裂重要性）迁移至深度模型，从而缩小深度模型与树模型之间的性能差距？
    \item 深度模型的注意力机制能否为欺诈检测提供有价值的可解释性？
    \item 在考虑实际业务成本的情况下，何种模型具有最佳的精度-成本权衡？
\end{enumerate}

\subsection{本文贡献}

本文的主要贡献如下：

\begin{enumerate}
    \item \textbf{统一评估框架下的异构架构横向比较}：在同一数据集和统一评估框架下，比较了两种异构深度表格架构（FT-Transformer的全局自注意力 vs. DeepFM的显式特征交互）与XGBoost基线的性能，揭示了不同归纳偏置在PCA匿名化特征空间中的表现差异；
    \item \textbf{树归纳偏置迁移方法}：提出了Tree-Supervised知识蒸馏方案，将XGBoost的软标签分布（方案3A）和特征分裂重要性先验（方案3B）迁移至FT-Transformer，验证了树模型的轴对齐归纳偏置可通过蒸馏增强深度模型的表格数据处理能力；
    \item \textbf{多维可解释性分析}：结合SHAP全局特征归因、个体瀑布图和FT-Transformer注意力权重，从多个视角解读模型决策逻辑，并诊断了自注意力在PCA空间中的均匀塌缩问题；
    \item \textbf{成本敏感评估}：引入业务成本矩阵（$C_{FP}=1$, $C_{FN}=10$），揭示AUPRC最优与业务成本最优之间的非单调关系，为实际部署提供帕累托前沿分析。
\end{enumerate}

\subsection{论文结构}

本文的其余部分组织如下：第二节回顾相关工作；第三节描述方法论；第四节呈现实验结果；第五节进行多维分析（可解释性、校准、成本）；第六节讨论商业启示、局限性和未来工作；第七节总结全文。
